\documentclass{article}
\usepackage{amsmath,amssymb,amsthm}
\usepackage[margin=1in]{geometry}

\newtheorem{theorem}{Theorem}
\newtheorem{lemma}[theorem]{Lemma}
\newtheorem{claim}[theorem]{Claim}
\newtheorem{remark}[theorem]{Remark}

\begin{document}

\title{A note on the proposed non-variation lemma for boundary strata}
\author{}
\date{}
\maketitle

\section{The statement as written is not provable}

The proposed assertion is the following.

\begin{claim}
Let \(V_0\subset \mathcal M_g\) be a complex smooth subvariety of complex dimension \(2\) which is totally geodesic for the Teichm\"uller metric. Let \(V_t\) be a holomorphic deformation satisfying the listed boundary assumptions. Fix a boundary stratum \(\Delta_\Gamma\), assumed to be a product
\[
\Delta_\Gamma \cong \mathcal M_{g_1,n_1}\times\cdots\times \mathcal M_{g_k,n_k},
\]
and put
\[
\partial V_t = \overline V_t\cap \Delta_\Gamma .
\]
Assume \(\dim \partial V_t=1\) and that, for some factor \(i\),
\[
\pi_i(\partial V_t)=\{Y_t\}\qquad\text{for all }t.
\]
Then \(Y_t=Y_0\) for all \(t\).
\end{claim}

The claim is false as stated. The boundary hypotheses (i) and (ii) do not force the holomorphic map
\[
\Delta\ni t\longmapsto Y_t\in \mathcal M_{g_i,n_i}
\]
to be constant. Totally geodesicity of the special fiber \(V_0\) is a condition only at \(t=0\); it does not prevent a holomorphic deformation of the subvariety from moving the boundary projection inside the interior of the factor \(\mathcal M_{g_i,n_i}\). In particular, the conclusion would require an additional hypothesis, e.g.\ that the whole family \(V_t\) remains totally geodesic, or that the projection \(Y_t\) is constant by definition.

To see explicitly that the stated boundary conditions alone cannot imply the conclusion, consider the standard plumbing model of a neighborhood of a boundary stratum. Near a point of \(\Delta_\Gamma\), the Deligne--Mumford compactification has holomorphic plumbing coordinates
\[
(q,y,c),
\]
where \(q\) is the smoothing parameter, \(y\) is a coordinate on the factor \(\mathcal M_{g_i,n_i}\), and \(c\) denotes the remaining product coordinates. The interior \(\mathcal M_g\) corresponds to \(q\neq 0\), while the boundary stratum \(\Delta_\Gamma\) is \(q=0\).

Let \(Y_t\) be any nonconstant holomorphic map from the unit disc into a small coordinate chart of \(\mathcal M_{g_i,n_i}\), with \(Y_0\) prescribed. For example, take \(Y_t=Y_0+t\) in a local chart. Let \(C\subset \prod_{j\neq i}\mathcal M_{g_j,n_j}\) be a fixed compact complex curve contained in the interior of the other factors, and choose a holomorphic function \(f\) on a neighborhood of \(C\). Define, locally,
\[
V_t=\{(q,y,c)\;:\; y=Y_t+q f(c),\ q\in D^*,\ c\in C\}.
\]
For each fixed \(t\), \(V_t\) is a smooth complex surface in \(\mathcal M_g\). The closure of \(V_t\) in the compactification meets \(\Delta_\Gamma\) in
\[
\partial V_t=\{(0,Y_t,c)\;:\;c\in C\}
=\{Y_t\}\times C.
\]
Hence
\[
\pi_i(\partial V_t)=\{Y_t\},
\]
which is nonconstant if \(Y_t\) is nonconstant. Thus the desired conclusion fails.

The family is smoothly trivial, because the map
\[
(q,c)\longmapsto (q,Y_t+q f(c),c)
\]
identifies the local models for different \(t\). Moreover, if \(C\) is chosen away from the boundary of the other factors, then no smaller stratum \(\Delta_{G'}\subset \partial\Delta_\Gamma\) is met, so condition (ii) is vacuous. Therefore all boundary compatibility assumptions are satisfied in this local model, while the conclusion is false.

This local model is realizable by the usual plumbing construction for degenerating families of stable curves. Consequently, the asserted lemma cannot be proved from the hypotheses listed in the question. A correct theorem would need an extra assumption, such as requiring the deformation \(V_t\) itself to remain totally geodesic, or requiring the projection map \(t\mapsto Y_t\) to be constant. Without such an additional hypothesis, the sought conclusion is not valid.

\end{document}